\documentclass[11pt]{article}
\usepackage[margin=0.8in]{geometry}
\usepackage{amsmath,amssymb,amsthm}
\usepackage{hyperref}
\emergencystretch=3em
\newtheorem{theorem}{Theorem}
\newtheorem{lemma}{Lemma}
\title{Subdivision cannot separate graph configuration homotopy types}
\author{Local formal research package for conjecture 00000004110}
\date{10 October 2026}
\begin{document}
\maketitle
\noindent\textbf{Result.} For every finite graph and every finite subdivision obtained by
inserting vertices on edges, the continuous ordered and unordered configuration
spaces are homeomorphic for every particle count. They are therefore homotopy
equivalent for every count, in particular at $n=2$. The source-bound refutation
targets its additional assertion that the four-claw graph is the smallest
example: no ordinary subdivision of that actual graph can be an example.

\section*{Original and scope}
The source defines subdivision as inserting vertices on edges. It then asserts
the existence of a graph and its proper subdivision whose configuration homologies
are isomorphic for all $n$, but whose configuration spaces are not homotopy
equivalent already for $n=2$, with the four-claw graph as the smallest example.
Both language versions express this same claim. The original does not explicitly
restrict all graphs to finite ones. Our general finite-graph theorem therefore
does not claim to settle an unrestricted infinite-graph existential by itself.
Instead, the smallest-example assertion necessarily requires a separation witness
whose original graph is the four-claw. We disprove that necessary witness clause.
No additional homology computation or minimality search is needed to refute the
whole conjunction. Arbitrary edgewise subdivisions of infinitely many original
edges are outside the formal general theorem's scope.

The formal graph presentation consists of finite vertex and edge types and
source/target maps. Vertices are discrete. Its realization is the quotient of the
disjoint union of its vertices and one closed real interval $[0,1]$ per edge,
identifying the interval's endpoints with the designated vertices. This is the
standard one-dimensional graph realization. Isolated vertices, loops and parallel
edges are allowed. Restricting to simple graphs therefore remains covered.
The result concerns continuous configurations on this realization, not selections
of combinatorial vertices or a discretized cubical approximation.

\section*{A genuine insertion homeomorphism}
For an edge $e$ from $u$ to $v$, insert a new vertex $w$ and a second edge. The
old edge label now ends at $w$, and the new edge runs from $w$ to $v$. All other
edges and incidences are preserved. On the realization, the collapse to the
original graph sends the first edge's coordinate $t$ to $t/2$, the second
edge's coordinate $t$ to $(1+t)/2$, and the new vertex to the original edge's
midpoint. Original vertices and unaffected edge coordinates remain fixed.

Its inverse sends an original edge point with $t\leq 1/2$ to the first edge at
$2t$, and with $t>1/2$ to the second edge at $2t-1$. At the dividing point the
two formulas agree in the quotient: the end of the first edge and the start
of the second both equal $w$. The formal implementation uses globally defined
clamped coordinates $\min(1,2t)$ and $\max(0,2t-1)$ to establish continuity
on the closed halves. The endpoint identifications are checked explicitly;
quotient lifting yields continuous maps. Direct coordinate calculations, including
the midpoint and both endpoints, prove both inverse identities. Thus the
homeomorphism is constructed rather than assumed.

\section*{Presentation identity and finite subdivisions}
Bijective relabelings of vertices and edges that preserve endpoints induce a
homeomorphism by applying the vertex bijection and carrying each interval, with
unchanged coordinate, to its new edge label. The inverse uses the inverse
bijections. Endpoint compatibility and continuity on each component are proved
before quotient lifting.

An edge's orientation is a presentation choice for an undirected graph. Reversing
selected edges swaps their source and target and replaces $t$ by $1-t$. This is
continuous, preserves the actual endpoint relation and is an involution. The
formal development further proves that an arbitrary bijection preserving each
edge's unordered endpoint pair factors through such reorientation and an
endpoint-preserving relabeling.

The inductive relation \texttt{Subdivision G H k} is the finite closure generated
by these actual presentation identities and one-edge insertions. Its index $k$
counts edge insertions: identity, relabeling and orientation changes contribute
zero; insertion contributes one; composition adds counts. A proper subdivision
has $k>0$. Its definition contains no homeomorphism or homotopy-equivalence
hypothesis. Induction composes the already constructed elementary homeomorphisms.
Any finite placement of new vertices in edge interiors gives precisely this
combinatorial operation; the chosen midpoint parametrization is harmless because
each resulting edge is a full closed interval with its endpoints glued.

\section*{Every ordinary four-claw subdivision}
The actual four-claw presentation has vertex type $\operatorname{Fin}(5)$,
edge type $\operatorname{Fin}(4)$, source constantly $0$ and target $i+1$ for
edge $i$. Thus it has one center, four distinct leaves and exactly their four
joining edges.

An ordinary subdivision replaces each edge by a nonempty finite path whose
interior vertices are fresh and independent from the other replacement paths.
This is the standard graph-theoretic convention: see Diestel, \emph{Graph Theory},
sixth edition, Sections 1.3 and 1.7, in the author's
\href{https://www.math.uni-hamburg.de/home/diestel/books/graph.theory/preview/Ch1.pdf}
{Chapter 1 preview}.
The four-claw has only four edges, so its ordinary subdivisions always require
finitely many insertions, irrespective of whether other graphs in the original
existential are allowed to be infinite.

The formal path-replacement operation adds $k$ interior vertices indexed by
$\operatorname{Fin}(k)$ and $k$ fresh edges. The old selected edge becomes the
first path edge. It joins the original source to interior vertex $0$ when $k>0$;
fresh edge $j$ joins interior vertex $j$ to $j+1$, or to the original target
at the final step. For $k=0$ the original edge stays intact. Induction on $k$
proves this actual path presentation belongs to the insertion closure. At the
induction step, insert one vertex on the first path edge; the explicit bijections
shift every old interior index and fresh-edge index by one, and give the new
vertex and new edge index zero. Every endpoint equation is checked.

Apply this replacement independently to all four original edges with arbitrary
counts $k_0,k_1,k_2,k_3$. The disjoint-sum label types make all added vertices
and edges fresh. Ordinary subdivision membership then allows actual vertex and
edge bijections preserving every unordered endpoint pair, hence arbitrary labels
and orientations. Its structural witness consists of the four path lengths and
these incidence bijections; it contains no homeomorphism, homotopy equivalence
or insertion-closure premise. The proved coverage theorem derives the insertion
closure, with count $k_0+k_1+k_2+k_3$, from that ordinary finite-path structure.
This covers the usual four independent finite-path replacement definition;
independent semantic review must verify this source correspondence.

\section*{Actual configuration spaces}
For a topological space $X$, define
\[
 F_n(X)=\{f:\{0,\ldots,n-1\}\longrightarrow X\mid f\text{ is injective}\},
\]
with the subspace topology inherited from the finite product $X^n$.
The unordered space $C_n(X)$ is the orbit quotient under permutations of the
indices, with its quotient topology. The formal orbit relation is the actual
existence of a permutation taking one configuration to the other.

If $h:X\longrightarrow Y$ is a homeomorphism, the coordinatewise map
$f\mapsto h\circ f$ preserves injectivity. Its inverse is
$g\mapsto h^{-1}\circ g$; both are continuous in the product/subspace topology.
It commutes with every index permutation, hence descends to a homeomorphism of
the unordered orbit quotients, again with a checked continuous inverse. This
construction works uniformly for every $n\in\mathbb N$, including zero.

\begin{theorem}
For finite graph presentations $G,H$ and any finite subdivision $H$ of $G$,
$F_n(|G|)$ and $F_n(|H|)$ are homeomorphic, as are $C_n(|G|)$ and $C_n(|H|)$,
for every $n\in\mathbb N$.
\end{theorem}
\begin{proof}
Compose the insertion and presentation homeomorphisms and apply the coordinatewise
and orbit-quotient constructions above.
\end{proof}

Mathlib's \texttt{ContinuousMap.HomotopyEquiv} consists of continuous maps in
both directions whose inverse composites are homotopic to the identities. Its
\texttt{Homeomorph.toHomotopyEquiv} supplies these witnesses from a homeomorphism:
the composites equal the identities, so constant interval homotopies suffice.
Applying it to the theorem's constructed homeomorphisms proves the actual
homotopy equivalences. Therefore the negated existential separation statements
at $n=2$ follow immediately for both conventions.

\section*{Lean correspondence and verification}
\begin{enumerate}
\item \texttt{ConnectedPartial.lean}: actual ordered and unordered configurations,
endpoint-cell realization, explicit insertion and inverse continuity.
\item \texttt{Relabeling.lean}: vertex/edge relabeling homeomorphism.
\item \texttt{Reorientation.lean}: genuine interval reversal and graph-cell identity.
\item \texttt{Subdivision.lean}: finite graph presentations, actual undirected
presentation identity, counted insertion closure and configuration homeomorphisms.
\item \texttt{Disproof.lean}: actual Mathlib homotopy equivalences and both
finite-graph separation contradictions.
\item \texttt{PathReplacement.lean}: concrete finite paths and their proved
membership in the insertion closure.
\item \texttt{FourClaw.lean}: actual $K_{1,4}$, ordinary four-path subdivision
coverage and both ordered and unordered four-claw separation contradictions.
\end{enumerate}
The package pins Lean 4.33.0 and every public dependency by Git commit. The graph
and configuration bridges are proved; no central bridge is an extra assumption.
Independent semantic review and the manager's deterministic build, type, axiom,
PDF and fresh-environment checks are required before local acceptance. This report
itself makes no claim of external submission or accepted solve status.
\end{document}
